% ID: SEL-D0041-OA-E02
% ESTADO: validada. S042, cotejada con examen y criterios.
\textbf{a)} El tramo racional es $2-13/(x+4)$.
Cuando $x\to-4^-$ tiende a $+\infty$ y cuando $x\to-4^+$
tiende a $-\infty$. La asíntota vertical es $\boxed{x=-4}$.
Cuando $x\to-\infty$ el tramo racional tiende a 2,
así que $y=2$ es la asíntota horizontal por la izquierda.
El tramo cúbico domina a cualquier recta al tender $x$
a $+\infty$. En $2$ ambos límites son finitos, de modo que
ese empalme tampoco produce una asíntota. El corte con el eje
vertical es $\boxed{(0,-5/4)}$.
\begin{center}
\setlength{\unitlength}{7mm}
\begin{picture}(12,7)(-7,-2)
\qbezier(-6.5,0)(-2,0)(2.5,0)\put(2.6,0){$x$}
\qbezier(0,-1.7)(0,1.2)(0,4.1)\put(0,4.2){$y$}
{\color{blue}\qbezier(-4,-1.7)(-4,1.2)(-4,4.1)
\qbezier(-6.5,2)(-2,2)(2.5,2)}
\put(-4,4.25){\makebox(0,0)[b]{$x=-4$}}
\put(2.7,2){$y=2$}
\put(0,-1.25){\makebox(0,0){$\bullet$}}
\put(0.2,-1.25){$(0,-5/4)$}
\end{picture}
\end{center}
\textbf{b)} $f(-3)=-11$ y $f'(x)=13/(x+4)^2$ en el tramo
racional. Así, $f'(-3)=13$ y la tangente
$y+11=13(x+3)$ queda $\boxed{y=13x+28}$.
